\documentclass{article}%
\usepackage[T1]{fontenc}%
\usepackage[utf8]{inputenc}%
\usepackage{lmodern}%
\usepackage{textcomp}%
\usepackage{lastpage}%
\usepackage{geometry}%
\geometry{margin=1in,headheight=14pt,headsep=25pt}%
\usepackage{hyperref}%
\usepackage{enumitem}%
\usepackage{fancyhdr}%
\usepackage{xcolor}%
\usepackage{url}%
\usepackage{breakurl}%
%

            \hypersetup{
                colorlinks=true,
                linkcolor=blue,
                filecolor=magenta,
                urlcolor=blue
            }
            \pagestyle{fancy}
            \fancyhf{}
            \rhead{Generated on \today}
            \cfoot{\thepage}
            
            % Configure enumeration settings
            \setlist[enumerate]{label=\Alph*.}
            \setlist[enumerate]{itemsep=0.5em}
            
            % Define a command for red text
            \newcommand{\incorrect}[1]{\textcolor{red}{#1}}
        %
\lhead{Lecture: 2024{-}10{-}22{-}Convexity}%
\title{Practice Questions for 2024{-}10{-}22{-}Convexity}%
\author{Generated by Scribe.AI}%
\date{\today}%
%
\begin{document}%
\normalsize%
\maketitle%
\section{Questions}%
\label{sec:Questions}%
1. According to Definition 1, which of the following statements is true about a convex set S in $\mathbb{R}^m$?%
\begin{enumerate}[label=\Alph*.]%
\item For any two points $z_1$ and $z_2$ in S, the line segment connecting them, defined by $tz_1 + (1-t)z_2$ where $t \ge 0$, is entirely contained within S.%
\item For any two points $z_1$ and $z_2$ in S, the line segment connecting them, defined by $tz_1 + (1-t)z_2$ where $0 < t < 1$, is entirely contained within S.%
\item For any two points $z_1$ and $z_2$ in S, the line segment connecting them, defined by $tz_1 + (1-t)z_2$ where $t \le 0$, is entirely contained within S.%
\item For any two points $z_1$ and $z_2$ in S, the line segment connecting them, defined by $tz_1 + (1-t)z_2$ where $t = 1$, is entirely contained within S.%
\item For any two points $z_1$ and $z_2$ in S, the line segment connecting them, defined by $tz_1 + (1-t)z_2$ where $t$ is any real number, is entirely contained within S.%
\end{enumerate}%
\vspace{1em}%
2. According to Theorem 10.1, which of the following is a necessary and sufficient condition for a set S to be convex?%
\begin{enumerate}[label=\Alph*.]%
\item S contains at least one convex combination of its points.%
\item S contains all convex combinations of its points.%
\item S contains no convex combinations of its points.%
\item S contains only the endpoints of convex combinations of its points.%
\item S contains only a subset of convex combinations of its points.%
\end{enumerate}%
\vspace{1em}%
3. In the proof of Theorem 10.1, when demonstrating that if S is convex, then it contains all convex combinations of its points, how is the case for n=3 handled?%
\begin{enumerate}[label=\Alph*.]%
\item It is shown that a convex combination of three points can be expressed as a convex combination of two points.%
\item It is shown that a convex combination of three points can be expressed as a convex combination of a convex combination of two points and a third point.%
\item It is shown that a convex combination of three points can be expressed as a linear combination of the three points.%
\item It is shown that a convex combination of three points can be expressed as a convex combination of four points.%
\item It is shown that a convex combination of three points can be expressed as a single point.%
\end{enumerate}%
\vspace{1em}

%
\newpage%
\section{Answers}%
\label{sec:Answers}%
1. According to Definition 1, which of the following statements is true about a convex set S in $\mathbb{R}^m$?%
\begin{enumerate}[label=\Alph*.]%
\item \incorrect{Answer A is incorrect because it includes $t=0$ and $t=1$, which are the endpoints of the line segment, not the line segment itself.}%
\item Answer B is correct because it accurately reflects the definition of a convex set, where the line segment connecting any two points within the set must also be within the set, with $0 < t < 1$ defining the line segment.%
\item \incorrect{Answer C is incorrect because it uses $t \le 0$, which is not the correct range for the parameter $t$ in the definition of a convex set.}%
\item \incorrect{Answer D is incorrect because it only considers the case where $t=1$, which is just one of the endpoints of the line segment, not the line segment itself.}%
\item \incorrect{Answer E is incorrect because it allows $t$ to be any real number, which does not define the line segment between $z_1$ and $z_2$.}%
\end{enumerate}%
\vspace{1em}%
2. According to Theorem 10.1, which of the following is a necessary and sufficient condition for a set S to be convex?%
\begin{enumerate}[label=\Alph*.]%
\item \incorrect{Answer A is incorrect because the theorem requires that S contains *all* convex combinations, not just one.}%
\item Answer B is correct because Theorem 10.1 states that a set S is convex if and only if it contains all convex combinations of its points.%
\item \incorrect{Answer C is incorrect because a convex set must contain convex combinations of its points.}%
\item \incorrect{Answer D is incorrect because a convex set must contain the entire convex combination, not just the endpoints.}%
\item \incorrect{Answer E is incorrect because a convex set must contain *all* convex combinations, not just a subset.}%
\end{enumerate}%
\vspace{1em}%
3. In the proof of Theorem 10.1, when demonstrating that if S is convex, then it contains all convex combinations of its points, how is the case for n=3 handled?%
\begin{enumerate}[label=\Alph*.]%
\item \incorrect{Answer A is incorrect because the proof does not directly reduce the convex combination of three points to a convex combination of two points, but rather uses a nested approach.}%
\item Answer B is correct because the proof for n=3 rewrites the convex combination of three points as a convex combination of a convex combination of two points and the third point, leveraging the base case (n=2).%
\item \incorrect{Answer C is incorrect because the proof specifically uses convex combinations, not just any linear combination.}%
\item \incorrect{Answer D is incorrect because the proof does not involve a convex combination of four points.}%
\item \incorrect{Answer E is incorrect because a convex combination of three points is not generally a single point.}%
\end{enumerate}%
\vspace{1em}

%
\end{document}